% ID: SEL-D0082-OA-E02
% Tanda S082.
\textbf{a)} Para ser derivable en $2$, la función debe ser continua allí y sus derivadas laterales deben coincidir. La continuidad exige
\[
 4a+6=4-2b-4=-2b.
\]
Las derivadas de los tramos son $2ax+3$ y $2x-b$, por lo que la igualdad de pendientes da $4a+3=4-b$. Despejando $b=1-4a$ y sustituyendo en la continuidad,
\[
 4a+6=-2+8a\quad\Longrightarrow\quad
 \boxed{a=2,\ b=-7}.
\]
Con estos valores, los dos tramos se unen con valor 14 y pendiente 11, confirmando continuidad y derivabilidad.

\textbf{b)} El punto de tangencia de $g$ es $(0,-2)$. Por la regla del cociente,
\[
 g'(x)=\frac{x-1-(x+2)}{(x-1)^2}=-\frac3{(x-1)^2}.
\]
La pendiente en $0$ es $-3$. La tangente es $y+2=-3x$, es decir, $\boxed{y=-3x-2}$.
